\chapter{Programar con IA: el método detrás de REM}\label{cap:metodo}

\begin{objbox}
Aprenderás cómo se reparte el trabajo entre una persona y un asistente de programación, qué prácticas hicieron que REM saliera adelante y qué errores reales cometió la IA (y cómo se detectaron).
\end{objbox}

\section{Quién hace qué}
REM se construyó en conversación con un asistente de programación con IA (en este caso, Claude Code, que lee y escribe archivos, ejecuta comandos y pruebas en el repositorio). El reparto fue este:

\begin{center}
\begin{tabular}{@{}p{0.2\linewidth}p{0.36\linewidth}p{0.36\linewidth}@{}}
\toprule
 & \textbf{La persona} & \textbf{La IA} \\
\midrule
Producto & Decide qué se construye, para quién y por qué. & Propone alternativas y señala riesgos. \\
Diseño & Aprueba la dirección. & Redacta el plan y lo compara con opciones. \\
Código & Revisa y prueba. & Escribe, ejecuta pruebas, corrige. \\
Riesgo & Decide qué es aceptable. & Avisa de lo que no se ha podido verificar. \\
Publicación & Da el visto bueno. & Prepara el cambio; no sube sin permiso. \\
\bottomrule
\end{tabular}
\end{center}

\begin{ideabox}
Programar con IA no es «pedir una aplicación». Es dirigir un trabajo en pasos pequeños, con criterios de «hecho» comprobables y una persona que mantiene el juicio sobre lo que importa.
\end{ideabox}

\section{Las prácticas que funcionaron}

\subsection{Producto antes que código}
Antes de escribir la parte de clientes (REM Clientes) se redactó un \textbf{documento de producto}, \emph{sin código}, para discutir: la idea en una frase, por qué puede ganar y por qué puede fracasar, principios, seguridad, qué no se haría. Todo lo no comprobado quedó marcado con «(verificar)». Esa disciplina evitó construir un prototipo sobre suposiciones sin examinar.

\subsection{Una hoja de ruta ordenada por valor y esfuerzo}
El archivo \ruta{ROADMAP.md} lista lo hecho y lo siguiente, en orden. La regla escrita en su cabecera: \emph{«Una función cada vez, con pruebas y commit local»}. Cada función se cerraba antes de empezar la siguiente.

\subsection{Plan primero en lo grande}
Para cambios amplios (la aplicación de Mac, el canal con clientes) se pidió primero un \textbf{plan y un diseño}, y solo tras aprobarlo se empezó a implementar. Una pregunta de elección a tiempo ahorra horas: ¿aplicación nativa o web envuelta? ¿importar con IA o con reglas? ¿copias en disco externo o en la nube?

\subsection{Pruebas como red de seguridad}
La suite (\texttt{npm test}) tiene unas 270 pruebas sin dependencias, incluidas pruebas con servidores simulados de correo y de IA. Cada función nueva llegó con sus pruebas. Cuando la IA cambia algo, las pruebas dicen en segundos si rompió otra cosa.

\subsection{Verificar en el navegador, no solo en el código}
Un test verde no demuestra que la pantalla se vea bien. Se comprobó cada cambio de interfaz en un navegador real, con capturas, en escritorio y móvil, y en tema claro y oscuro. Una auditoría automática de botones táctiles pasó de unos 30 por debajo del tamaño recomendado a 3.

\subsection{Confirmar antes de publicar}
La IA hacía commits \emph{locales}; subir al repositorio público requería una confirmación explícita de la persona cada vez. Es una barrera sencilla y muy eficaz.

\subsection{Datos ficticios siempre}
Todo lo que se enseña (capturas, vídeos, ejemplos, este libro) usa datos inventados y marcados como tales: NIF de ceros, correos en \texttt{.invalid}, empresas imaginarias.

\section{Los errores reales de la IA (y cómo se encontraron)}
Una lección honesta: el asistente se equivoca. Lo importante es tener mecanismos que lo detecten.

\begin{center}
\small
\begin{tabular}{@{}p{0.34\linewidth}p{0.3\linewidth}p{0.28\linewidth}@{}}
\toprule
\textbf{Error} & \textbf{Cómo apareció} & \textbf{Qué lo detectó} \\
\midrule
El relé nunca borraba los mensajes entregados: tras unos miles, bloqueaba a la empresa. & Con la cola llena respondía «503». & Una \textbf{prueba de carga} con 1\,000 clientes. \\
Al vaciarse la cola, los identificadores se reutilizaban y la empresa ignoraba mensajes nuevos. & Mensajes que nunca llegaban. & Un test de regresión escrito tras el fallo. \\
En la lectura de la web con el modelo local se usaba mal el resultado de la función. & La IA local «no respondía». & Un test con un servidor simulado de IA. \\
Los botones «Pago» y «Estado» quedaban fuera de la cabecera de la conversación. & Se veían como barras a todo ancho. & Al \textbf{hacer capturas} para el vídeo. \\
La validación de NIF aceptaba un CIF con dígito de control incorrecto y un valor vacío. & Un NIF inventado pasaba. & Un test escrito para el nuevo extractor. \\
Valores de ejemplo con un dominio de correo que podría ser real. & Al revisar las capturas. & La \textbf{revisión de privacidad} antes de publicar. \\
\bottomrule
\end{tabular}
\end{center}

\begin{warnbox}
La IA tiende a explicar con seguridad lo que no ha comprobado. La defensa no es desconfiar de todo, sino pedir siempre la \emph{evidencia}: una prueba que pasa, una captura, un archivo que se abre, un número que se puede contar.
\end{warnbox}

\section{Cómo pedir las cosas}
Un buen encargo tiene cuatro partes: \textbf{contexto}, \textbf{objetivo}, \textbf{restricciones} y \textbf{criterio de hecho}.

\begin{center}
\small
\begin{tabular}{@{}p{0.45\linewidth}p{0.45\linewidth}@{}}
\toprule
\textbf{Encargo flojo} & \textbf{Encargo claro} \\
\midrule
«Haz las copias de seguridad.» & «Copias cifradas con AES-256-GCM en una carpeta que elija la persona; una al día y otra al salir; que el asistente pruebe una restauración antes de darlas por buenas. Con pruebas.» \\
«Mejora la seguridad.» & «Enumera qué puede hacer un correo malicioso en esta bandeja y qué barrera lo impide; añade pruebas para cada caso.» \\
«Que la IA no haga cosas peligrosas.» & «Las acciones de enviar, emitir y borrar tienen un suelo en código que ninguna configuración puede bajar; muéstrame la prueba.» \\
\bottomrule
\end{tabular}
\end{center}

\section{Lo que no se delega}
\begin{itemize}
\item \textbf{Qué riesgos son aceptables.} La IA los describe; la persona decide.
\item \textbf{Lo que se afirma en público.} Cualquier mensaje o vídeo que vaya a terceros se revisa para no prometer lo que no está verificado (por ejemplo, llamar «seguro» al prototipo del canal con clientes).
\item \textbf{La licencia y la autoría.} Son decisiones del titular.
\end{itemize}

\section{Lista de revisión para código generado por IA}
\begin{enumerate}
\item ¿Hay una prueba que falle sin el cambio y pase con él?
\item ¿Qué ocurre con entradas raras (vacías, enormes, con caracteres especiales, de otro usuario)?
\item ¿Hay un camino en el que la IA pueda actuar sin aprobación humana?
\item ¿Se han añadido dependencias? ¿Hacían falta?
\item ¿Se ve bien en pantalla pequeña y con tema oscuro?
\item ¿Queda constancia de lo que no se ha podido verificar?
\end{enumerate}

\begin{checkbox}
Elige una función pequeña de REM (por ejemplo, «archivar un cliente»). Escribe un encargo con contexto, objetivo, restricciones y criterio de hecho; pídeselo a un asistente de programación; y evalúa el resultado con la lista de revisión anterior.
\end{checkbox}
